\section{Summary}
\begin{frame}{}
    \LARGE Learning from Videos: \textbf{Summary}
\end{frame}

\begin{frame}{Key Takeaways}
    \large
    \begin{itemize}
        \item Video understanding evolves from single-frame analysis to capturing global context using convolution, recurrence, and attention mechanisms.
        \item Temporal modeling is essential for tasks like action recognition and future prediction.
        \item Self-supervised learning reduces the need for extensive manual annotation.
    \end{itemize}
\end{frame}

\begin{frame}{Limitations and Open Problems}
    \begin{itemize}
        \item \textbf{Computational cost:} 3D convolutions are expensive (C3D: 39.5 GFLOP per clip), and full spatio-temporal attention grows quadratically with sequence length.
        \item \textbf{Short clips:} Most models see only 16--64 frames at a time; structure spanning minutes still relies on sparse sampling, recurrence, or external memory.
        \item \textbf{Annotation cost:} Frame-level labels are expensive, and self-supervised pretext tasks do not transfer equally well to every downstream task.
        \item \textbf{Appearance shortcuts:} Models can classify many actions from scene appearance alone while ignoring motion, which is why motion is evaluated separately.
        \item \textbf{Open problems:} Efficient architectures for long videos, learning from egocentric and streaming video, and scaling self-supervision to unlabeled video at internet scale.
    \end{itemize}
\end{frame}